\subsection{Closed points over a finite field: proof of Corollary~\ref{cor:closed}}
Let $k$ be finite, $X$ a hyperbolic curve over $k$, and suppose that $\tilde x\ne\tilde x'$ are closed points of $\widetilde X$ such that $D_0=D_{\tilde x}\cap D_{\tilde x'}$ is open in $D_{\tilde x}$ and in $D_{\tilde x'}$. By \cite[Proposition 1.8(ii)]{ST09} neither point lies above a point at infinity, and by \cite[Proposition 1.8(iii)]{ST09} the images $\bar x$ and $\bar x'$ of the two points in $X_{\overline k}$ are distinct. The decomposition group of a closed point which is not at infinity meets $\Delta_X$ trivially, so the projection to $G_k$ is injective on $D_{\tilde x}$ and on $D_{\tilde x'}$, with open image. Hence the image of $D_0$ is $G_L$ for a finite extension $L$ of $k$ in $\overline k$.

The inverse image $\Pi_{X_L}$ of $G_L$ in $\Pi_X$ is the geometrically prime-to-$p$ fundamental group of $X_L=X\otimes_kL$, with the same pro-cover. The group $D_{\tilde x}\cap\Pi_{X_L}$ is the decomposition group of $\tilde x$ in $\Pi_{X_L}$. It contains $D_0$ and maps injectively to $G_L$, while $D_0$ maps onto $G_L$; so it equals $D_0$. Hence the image $x_L$ of $\tilde x$ in $X_L$ is an $L$-rational point, and $D_0$ is the image of a section $s_{x_L}$ of $\Pi_{X_L}\to G_L$. The same holds for $\tilde x'$ and its image $x'_L$. The sections $s_{x_L}$ and $s_{x'_L}$ have the same image $D_0$, so they are equal, and Theorem~\ref{thm:main}, applied to $X_L$ over $L$, gives $x_L=x'_L$. But an $L$-rational point of $X_L$ has exactly one point of $X_{\overline k}$ above it, so $\bar x=\bar x'$, a contradiction. This proves Corollary~\ref{cor:closed}.

\section{What remains open}
For the geometrically pro-$\Sigma$ quotient, with $\Sigma$ a set of primes smaller than the set of all primes different from $p$, neither Theorem~\ref{thm:main} nor Corollary~\ref{cor:closed} is proved. For $\Sigma$ cofinite a finite exceptional set is available \cite[Proposition 1.8(vi)]{ST09}, which is what Step~4 needs, and we have not checked whether the other steps go through; for $\Sigma$ not cofinite even the finiteness of the exceptional sets is open \cite[Introduction, pp.~4--5]{ST18}. For an infinite finitely generated field $K$ we have not treated the decomposition groups of closed points of higher degree; the argument of Corollary~\ref{cor:closed} is the natural starting point. Nothing is proved here about surjectivity, or about varieties of higher dimension.

\section*{Note on method}
The theorem and its proof were produced by GPT-6 Astra (ChatGPT Pro mode). Claude Opus 5.5 at effort max re-proved Steps~1 and~2 independently from the statements. In a separate check, Claude Opus 5.5 at effort max checked Steps~3 to~5 and Section~\ref{sec:open} line by line, corrected the genus-one argument (the lifted points are paired by a deck transformation; no translation is needed), wrote out the divisibility argument for the generalized Jacobian, and wrote the first version of this text. Claude Fable 5.2 at effort max edited the text into its present form and followed every step again. That reading made three repairs of wording in the proofs, none of which changes a statement: the extension $K'$ in the genus-one case must also split the fibre above $P$; Lemma 1.7 of \cite{ST09} is cited for what it states; and a superfluous appeal to cohomological dimension in Step~5 is removed. It also added the reformulation by covers after Theorem~\ref{thm:main}, the remark on elliptic curves, and Corollary~\ref{cor:closed} with its proof. Corollary~\ref{cor:closed} has not been checked by a second reader.

Of the sources, the passages of \cite{ST09} cited here were read at the source for this version: the Introduction, Remark 1.2, Definitions 1.3, 1.5, 1.6 and 1.10, the statements of Lemma 1.7, Proposition 1.8 and Corollary 1.11, and in \S 2 the statement of Theorem 2.2 and the passages on the fundamental extension. The Introduction of \cite{ST18} and Grothendieck's manuscript pages, in our transcription and translation, were read by Claude Opus 5.5 at effort max and not again for this version. The locators in \cite{SGA1} and \cite{WeilI} are those of the proof as it was produced; they were not reopened. \cite{Moch} is cited through \cite[\S 2]{ST09} and was not opened, and its journal reference is given from memory. \cite{Stix} could not be consulted.

No human has checked this note. Our text is dedicated to the public domain under CC0~1.0. The workbench is \expandafter\url\expandafter{\workbenchurl}; the permanent records are at \expandafter\url\expandafter{\recordurl}.

\begin{thebibliography}{9}
\bibitem{Archive} A.~Grothendieck, manuscripts, Archives Grothendieck, Universit\'e de Montpellier, cote 149, ms.~pp.~17--18 (scans 35--36 of the digitized cote). \url{https://grothendieck.umontpellier.fr/archives-grothendieck/}.
\bibitem{SGA1} A.~Grothendieck, \emph{Rev\^etements \'etales et groupe fondamental} (SGA~1), Lecture Notes in Math.\ 224, Springer, 1971: III 7.3--7.4; IX 4.10; X 3.8--3.9; XIII 2.9--2.12 and 5.2.
\bibitem{Moch} S.~Mochizuki, \emph{Absolute anabelian cuspidalizations of proper hyperbolic curves}, J.\ Math.\ Kyoto Univ.\ \textbf{47} (2007), 451--539: Propositions 1.6 and 1.8, Theorem 1.16, as cited in \cite[\S 2]{ST09}.
\bibitem{ST09} M.~Sa\"idi and A.~Tamagawa, \emph{A prime-to-$p$ version of Grothendieck's anabelian conjecture for hyperbolic curves over finite fields of characteristic $p>0$}, Publ.\ RIMS \textbf{45} (2009), 135--186; arXiv:0801.1452: Introduction; Remark 1.2; Definitions 1.3, 1.5, 1.10; Lemma 1.7; Proposition 1.8; Corollary 1.11; \S 2.
\bibitem{ST18} M.~Sa\"idi and A.~Tamagawa, \emph{A refined version of Grothendieck's anabelian conjecture for hyperbolic curves over finite fields}, J.\ Algebraic Geom.\ \textbf{27} (2018), 383--448; arXiv:1708.08773, Introduction, pp.~4--5.
\bibitem{Stix} J.~Stix, \emph{Rational points and arithmetic of fundamental groups: evidence for the section conjecture}, Lecture Notes in Math.\ 2054, Springer, 2013 (not consulted).
\bibitem{WeilI} P.~Deligne, \emph{La conjecture de Weil.~I}, Publ.\ Math.\ IH\'ES \textbf{43} (1974), 273--307: (1.5.1) and Th\'eor\`eme 1.6.
\end{thebibliography}

\end{document}
